\documentclass{article}
\usepackage[margin=1in]{geometry}
\usepackage{amsmath,amsfonts,mathtools,amsthm,amssymb}
\usepackage[l2tabu,orthodox]{nag}
\usepackage{microtype} % fixes spacing or whatever
\usepackage{enumitem}
\usepackage{tikz-cd}
\usepackage{xcolor}
\usepackage{physics}
\usepackage{mathrsfs}
\usepackage{cancel}
% \usepackage{libertine,libertinust1math}
\usepackage{eucal}
\usepackage[toc,page]{appendix}
\usepackage{pgfplots}
\pgfplotsset{compat=1.18}

\newtheorem{proposition}{Proposition}
\newtheorem{theorem}{Theorem}
\newtheorem{lemma}{Lemma}
\newtheorem{corollary}{Corollary}

\theoremstyle{definition}
\newtheorem{remark}{Remark}
\newtheorem{definition}{Definition}
\newtheorem{example}{Example}

% \usepackage[sc,osf]{mathpazo}   % With old-style figures and real smallcaps.
% \linespread{1.025}              % Palatino leads a little more leading
% % Euler for math and numbers
% \usepackage[euler-digits,small]{eulervm}
% \AtBeginDocument{\renewcommand{\hbar}{\hslash}}

\usepackage{etoolbox}
\usepackage{dynkin-diagrams}
\def\row#1/#2!{\dynkin{#1}{#2}\\}
\newcommand{\tble}[1]{
   \renewcommand*\do[1]{\row##1!}
   \[
      \begin{array}{ll}\docsvlist{#1}\end{array}
   \]
}

\allowdisplaybreaks

\renewcommand{\epsilon}{\varepsilon}
% \renewcommand{\phi}{\varphi}

\DeclareMathAlphabet{\mathpzc}{OT1}{pzc}{m}{it}

\newcommand{\goth}[1]{\mathfrak{#1}}
\newcommand{\red}[1]{#1_{\text{red}}}
\newcommand{\ad}[1]{#1_{\text{Ad}}}
\newcommand{\reg}[1]{#1_{\text{reg}}}
\newcommand{\sh}[1]{#1^{\text{sh}}}
\newcommand{\cpdv}[2]{\cfrac{\partial #1}{\partial #2}}

\newcommand{\fA}{\mathcal{A}}
\newcommand{\fB}{\mathcal{B}}
\newcommand{\fC}{\mathcal{C}}
\newcommand{\fD}{\mathcal{D}}
\newcommand{\fE}{\mathcal{E}}
\newcommand{\fF}{\mathcal{F}}
\newcommand{\fG}{\mathcal{G}}
\newcommand{\fH}{\mathcal{H}}
\newcommand{\fI}{\mathcal{I}}
\newcommand{\fJ}{\mathcal{J}}
\newcommand{\fK}{\mathcal{K}}
\newcommand{\fL}{\mathcal{L}}
\newcommand{\fM}{\mathcal{M}}
\newcommand{\fN}{\mathcal{N}}
\newcommand{\fO}{\mathcal{O}}
\newcommand{\fP}{\mathcal{P}}
\newcommand{\fQ}{\mathcal{Q}}
\newcommand{\fR}{\mathcal{R}}
\newcommand{\fS}{\mathcal{S}}
\newcommand{\fT}{\mathcal{T}}
\newcommand{\fX}{\mathcal{X}}
\newcommand{\fY}{\mathcal{Y}}
\newcommand{\fZ}{\mathcal{Z}}

\newcommand{\be}{\mathbf{e}}
\newcommand{\bx}{\mathbf{x}}
\newcommand{\bone}{\mathbf{1}}
\newcommand{\cx}{\bar{x}}
\newcommand{\cy}{\bar{y}}
\newcommand{\vx}{\vec{x}}
\newcommand{\vsigma}{\vec{\sigma}}
\newcommand{\tzeta}{\tilde{\zeta}}

\newcommand{\gm}{\goth{m}}
\newcommand{\gp}{\goth{p}}
\newcommand{\gF}{\goth{F}}
\newcommand{\gU}{\goth{U}}
\newcommand{\gV}{\goth{V}}

\newcommand{\A}{\mathbf{A}}
\newcommand{\C}{\mathbf{C}}
\newcommand{\F}{\mathbf{F}}
\newcommand{\G}{\mathbf{G}}
\newcommand{\bH}{\mathbf{H}}
\newcommand{\N}{\mathbf{N}}
\newcommand{\PP}{\mathbf{P}}
\newcommand{\Q}{\mathbf{Q}}
\newcommand{\R}{\mathbf{R}}
\newcommand{\Z}{\mathbf{Z}}

\newcommand{\dlog}{\dd\log}

\newcommand\srestr[2]{{\left.\kern-\nulldelimiterspace #1\vphantom{\small|} \right|_{#2}}}
\newcommand\restr[2]{{\left.\kern-\nulldelimiterspace #1 \vphantom{\big|} \right|_{#2}}}
\newcommand\gsec[2]{\Gamma({#1}, {#2})}
\newcommand\gsecx[1]{\Gamma(X, {#1})}

\DeclareMathOperator{\chr}{char}
\DeclareMathOperator{\rProj}{\mathbf{Proj}}
\DeclareMathOperator{\rSpec}{\mathbf{Spec}}
\DeclareMathOperator{\rHom}{\mathbf{Hom}}
\DeclareMathOperator{\rExt}{\mathbf{Ext}}
\DeclareMathOperator{\id}{id}
\DeclareMathOperator{\Frac}{Frac}
\DeclareMathOperator{\rk}{rank}
\DeclareMathOperator{\deter}{det}
\DeclareMathOperator{\pic}{Pic}
\DeclareMathOperator{\cacl}{CaCl}
\DeclareMathOperator{\trd}{tr.d.}
\DeclareMathOperator{\cl}{Cl}
\DeclareMathOperator{\depth}{depth}
\DeclareMathOperator{\codim}{codim}
\DeclareMathOperator{\Div}{Div}
\DeclareMathOperator{\coker}{coker}
\DeclareMathOperator{\len}{length}
\DeclareMathOperator{\height}{height}
\DeclareMathOperator{\supp}{Supp}
\DeclareMathOperator{\proj}{Proj}
\DeclareMathOperator{\im}{im}
\DeclareMathOperator{\Hom}{Hom}
\DeclareMathOperator{\Der}{Der}
\DeclareMathOperator{\spec}{Spec}
\DeclareMathOperator{\Aut}{Aut}
\DeclareMathOperator{\ch}{char}
\DeclareMathOperator{\tor}{Tor}
\DeclareMathOperator{\Ann}{Ann}
\DeclareMathOperator{\Syl}{Syl}
\DeclareMathOperator{\Sym}{Sym}
\DeclareMathOperator{\GL}{GL}
\DeclareMathOperator{\SL}{SL}
\DeclareMathOperator{\Stab}{Stab}
\DeclareMathOperator{\Perm}{Perm}
\DeclareMathOperator{\Orb}{Orb}
\DeclareMathOperator{\Gal}{Gal}
\DeclareMathOperator{\Supp}{Supp}
\DeclareMathOperator{\Ext}{Ext}
\DeclareMathOperator{\Tor}{Tor}
\DeclareMathOperator{\PGL}{PGL}
\DeclareMathOperator{\hd}{hd}
\DeclareMathOperator{\pd}{pd}
\DeclareMathOperator{\SO}{SO}
\DeclareMathOperator{\SU}{SU}
\DeclareMathOperator{\Sp}{Sp}
\DeclareMathOperator{\Oth}{O}
\DeclareMathOperator{\Uni}{U}
\DeclareMathOperator{\evf}{ev}
\DeclareMathOperator{\cH}{H}
\DeclareMathOperator{\dR}{R}
\DeclareMathOperator{\End}{End}
\DeclareMathOperator{\Hasse}{Hasse}
\DeclareMathOperator{\Num}{Num}

\author{James Lee}

% Solarized
% \pagecolor[RGB]{0,20,26}
% \color[RGB]{191,191,191}

% Sephia
% \pagecolor[RGB]{249,239,220}

% Grey
% \pagecolor[RGB]{200, 200, 200}

% Black
% \pagecolor[RGB]{0,0,0}
% \color[RGB]{255,255,255}

\title{Chapter 4, Section 2}

\begin{document}
\maketitle
\begin{enumerate} [label=\textbf{\arabic*.}, leftmargin=0em]

\item Use (2.5.3) to show that $\PP^n$ is simply connected.

\begin{proof}
  Let $f : X \to \PP^n$ be an étale covering.
  We may assume $X$ is connected.
  Then $X$ is smooth over $k$ since $f$ is étale, and $X$ is projective over $k$ since $f$ has finite, so $X$ has dimension $n$.

\end{proof}

\item \textit{Classification of Curves of Genus 2.} Fix an algebraically closed field $k$ of characteristic $\neq 2$.
\begin{enumerate} [label=(\alph*)]
  \item If $X$ is a curve of genus 2 over $k$, the canonical linear system $|K|$ determines a finite morphism $f : X \to \PP^1$ of degree 2 (Ex. 1.7). Show that it is ramified at exactly 6 points, with ramification index 2 at each one.
  Note that $f$ is uniquely determined, up to an automorphism of $\PP^1$, so $X$ determines an (unordered) set of 6 points of $\PP^1$, up to an automorphism of $\PP^1$.

  \item Conversely, given six distinct elements $\alpha_1, \dots, \alpha_6 \in k$, let $K$ be the extension of $k(x)$ determined by the equation $z^2 = (x - \alpha_1) \cdots (x-\alpha_6)$.
  Let $f : X \to \PP^1$ be the corresponding morphism of curves.
  Show that $g(X) = 2$, the map $f$ is the same as the one determined by the canonical linear system, and $f$ is ramified over the six points $x = \alpha_i$ of $\PP^1$, and nowhere else.

  \item Using (I, Ex. 6.6), show that if $P_1, P_2, P_3$ are three distinct points of $\PP^1$, then there exists a unique $\varphi \in \Aut{\PP^1}$ such that $\varphi(P_1) = 0$, $\varphi(P_2) = 1$, $\varphi(P_3) = \infty$.
  Thus in (a), if we order the six points of $\PP^1$, and then normalize by sending the first three to $0, 1, \infty$, respectively, we may assume that $X$ is ramified over $0, 1, \infty, \beta_1, \beta_2, \beta_3$, where $\beta_1, \beta_2, \beta_3$ are three distinct elements of $k$, $\neq 0, 1$.

  \item Let $\Sigma_6$ be the symmetric group on 6 letters.
  Define an action of $\Sigma_6$ on sets of three distinct elements $\beta_1, \beta_2, \beta_3$ of $k$, $\neq 0, 1$, as follows: reorder the set $0, 1, \infty, \beta_1, \beta_2, \beta_3$ according to a given element $\sigma \in \Sigma_6$, then renormalize as in (c) so that the three become $0, 1, \infty$ again.
  Then the last three are the new $\beta_1', \beta_2', \beta_3'$.

  \item Summing up, conclude that there is a one-to-one correspondence between the set of isomorphism classes of curves of genus 2 over $k$, and triples of distinct elements $\beta_1, \beta_2, \beta_3$ of $k$, $\neq 0, 1$, modulo the action of $\Sigma_6$ described in (d).
  In particular, there are many non-isomorphic curves of genus 2.
  We say that curves of genus 2 depend on three parameters, since they correspond to the points of an open subset of $\A_k^3$ modulo a finite group.
\end{enumerate}

\begin{proof} $ $ \vspace{0pt}
\begin{enumerate} [leftmargin=0cm]
  \item By Hurtwitz theorem, the degree of the ramification divisor is computed to be
  \begin{equation*}
    \deg{R} = 2g(X) - 2 - n(2g(\PP^1) - 2) = 6.
  \end{equation*}
  The ramification index at a point cannot exceed the degree of the morphism, so $f$ must be ramified at exactly 6 distinct points each having ramification index 2.

  \item Since $K$ has degree 2 over $k(x)$, the degree of $f$ is 2.

  \item Let $x$ be an affine coordinate on $\PP^1$ such that $P_i \neq \infty$ and identify $P_i$ with some $\alpha_i \in k$.
  The following automorphism will do
  \begin{equation*}
    \varphi(x) = \frac{\alpha_2 - \alpha_3}{\alpha_2 - \alpha_1}\frac{x - \alpha_1}{x - \alpha_3}.
  \end{equation*}

  \item[(e)] We want to show two curves of genus 2 are isomorphic if and only if they correspond to the same triples of distinct elements in $k$, $\neq 0, 1$, described in (d).
\end{enumerate} 
\end{proof}

\item[\textbf{5.}] \textit{Automorphisms of a Curve of Genus $\geq 2$.} Prove the theorem of Hurtwitz [1] that a curve $X$ of genus $g \geq 2$ over a field of characteristic 0 has at most $84(g - 1)$ automorphisms.
We will see later (Ex. 5.2) or (V, Ex. 1.11) that the group $G = \Aut{X}$ is finite.
So let $G$ have order $n$.
Then $G$ acts on the function field $K(X)$.
Let $L$ be the fixed field.
Then the field extension $L \subseteq K(X)$ corresponds to a finite morphism of curves $f : X \to Y$ degree $n$.
\begin{itemize}
  \item[(a)] If $P \in X$ is a ramification point, and $e_p = r$, show that $f^{-1}f(P)$ consists of exactly $n/r$ points, each having ramification index $r$.
  Let $P_1, \dots, P_s$ be a maximal set of ramification points of $X$ lying over distinct points of $Y$, and let $e_{P_i} = r_i$.
  Then show that Hurtwitz's theorem implies that
  \begin{equation*}
    (2g - 2)/n = 2g(Y) - 2 + \sum_{i = 1}^s (1 - 1/r_i).
  \end{equation*}
  \item[(b)] Since $g \geq 2$, the left hand side of the equation is $> 0$.
  Show that if $g(Y) \geq 0$, $s \geq 0$, $r_i \geq 2$, $i = 1, \dots, s$ are integers such that
  \begin{equation*}
    2g(Y) - 2 + \sum_{i = 1}^s (1 - 1 / r_i) > 0,
  \end{equation*}
  then the minimum value of this expression is $1 / 42$.
  Conclude that $n \geq 84(g - 1)$.
\end{itemize}

\begin{proof} $ $ \vspace{0pt}
\begin{enumerate} [leftmargin=0cm]
  \item[(a)]

  \item[(b)] The following values give the minimum value
  \begin{equation*}
    g(Y) = 0, \quad s = 3, \quad r_1 = \frac{1}{2}, \quad r_2 = \frac{1}{3}, \quad r_3 = \frac{1}{7}.
  \end{equation*}
  We conclude that
  \begin{equation*}
    n \leq 42(2g - 2) = 84(g - 1).
  \end{equation*}
\end{enumerate} 
\end{proof}

\newpage

\item[\textbf{6.}] \textit{$f_*$ for Divisors.} Let $f : X \to Y$ be a finite morphism of curves of degree $n$.
We define a homomorphism $f_* : \Div{X} \to \Div{Y}$ by $f_*(\sum n_i P_i) = \sum n_i f(P_i)$ for any divisor $D = \sum n_i P_i$ on $X$.
\begin{itemize}
  \item[(a)] For any locally free sheaf $\fE$ on $Y$, of rank $r$, we define $\det \fE = \wedge^r \fE \in \pic{Y}$ (II, Ex. 6.11).
  In particular, for any invertible sheaf $\fM$ on $X$, $f_*\fM$ is locally free of rank $n$ on $Y$, so we can consider $\det f_*\fM \in \pic{Y}$.
  Show that for any divisor $D$ on $X$,
  \begin{equation*}
    \det(f_*\fO_X(D)) \cong (\det f_*\fO_X) \otimes \fO_Y(f_*D).
  \end{equation*}
  Note in particular that $\det(f_*\fO_X(D)) \neq \fO_Y(f_*D)$ in general!.
  \item[(b)] Conclude that $f_*D$ depends only on the linear equivalence class of $D$, so there is an induced homomorphism $f_* : \pic{X} \to \pic{Y}$.
  Show that $f_* \circ f^* : \pic{Y} \to \pic{Y}$ is just multiplication by $n$.
  \item[(c)] Use duality for a finite flat morphisms (III, Ex. 6.10) and (III, Ex. 7.2) to show that
  \begin{equation*}
    \det f_*\Omega_X \cong (\det f_*\fO_X )^{-1} \otimes \Omega_Y^{\otimes n}.
  \end{equation*}
  \item[(d)] Now assume that $f$ is separable, so we have the ramification divisor $R$.
  We define the \textit{branch divisor} $B$ to be the divisor $f_*R$ on $Y$.
  Show that
  \begin{equation*}
    (\det f_*\fO_X)^2 \cong \fO_Y(-B).
  \end{equation*}
\end{itemize}

\begin{proof} $ $ \vspace{0pt}
\begin{itemize} [leftmargin=0cm]
  \item[(a)] Consider the case when $D$ is an effective divisor.
  The pushfoward $f_*$ is exact for finite morphisms, so applying $f_*$ to the exact sequence
  \[ \begin{tikzcd}
    0 \arrow[r] & \fO_X \arrow[r] & \fO_X(D) \arrow[r] & \fO_D \arrow[r] & 0
  \end{tikzcd}\]
  gives us another exact sequence
  \[ \begin{tikzcd}
    0 \arrow[r] & f_*\fO_X \arrow[r] & f_*\fO_X(D) \arrow[r] & f_*\fO_D \arrow[r] & 0.
  \end{tikzcd}\]
  Taking determinants, we have
  \begin{equation*}
    \det(f_*\fO_X(D)) \cong (\det f_*\fO_X) \otimes \det(f_*\fO_D)
  \end{equation*}

  \item[(b)] It suffices to show $f_*f^*Q = nQ$ for every point $Q \in Y$.
  By definition, $f_*f^*Q = n'Q$ for some $n' \geq 0$.
  On the otherhand, $\deg f_*D = \deg{D}$, and $\deg f^*D = n \deg{D}$ (II, 6.9), so $\deg f_*f^*D = n \deg{D}$.
  Hence, $n' = n$.

  \item[(c)]

  \item[(d)] Let $K_X$ and $K_Y$ be canonical divisors of $X$ and $Y$, respectively.
  They satisfy the following relation (2.3)
  \begin{equation*}
    f^*K_Y - K_X \sim -R.
  \end{equation*}
  Applying $f_*$ gives
  \begin{equation*}
    nK_Y - f_*K_X \sim - B,
  \end{equation*}
  where we used the fact that $f_*f^*K_Y = nK_Y$ (b).
  On the otherhand, the results of (a) applied to $D = K_X$ and (c) give
  \begin{equation*}
    \fO_Y(nK_Y) \cong (\det f_*\fO_X) \otimes (\det f_*\Omega_X), \quad \fO_Y(-f_*K_X) \cong (\det f_*\Omega_X) \otimes (\det f_*\Omega_X)^{-1},
  \end{equation*}
  so combining everything, we get
  \begin{equation*}
    (\det f_*\fO_X)^2 \cong \fO_Y(nK_Y - f_*K_X) \cong \fO_Y(-B).
  \end{equation*}
\end{itemize} 
\end{proof}

\item[\textbf{7.}] \textit{Étale Covers of Degree 2.} Let $Y$ be a curve over a field $k$ of characteristic $\neq 2$.
We show there is a one-to-one correspondence between finite étale morphisms $f : X \to Y$ of degree 2, and 2-torsion elements of $\pic{Y}$, i.e., invertible sheaves $\fL$ on $Y$ with $\fL^2 \cong \fO_Y$.
\begin{itemize}
  \item[(a)] Given an étale morphism $f : X \to Y$ of degree 2, there is a natural map $\fO_Y \to f_*\fO_X$.
  Let $\fL$ be the cokernel.
  Then $\fL$ is an invertible sheaf on $Y$, $\fL \cong \det f_*\fO_X$, and so $\fL^2 \cong \fO_X$ by (Ex. 2.6).
  Thus an étale cover of degree 2 determines a 2-torsion element in $\pic{Y}$.
  \item[(b)] Conversely, given a 2-torsion element $\fL$ in $\pic{Y}$, define an $\fO_Y$-algebra structure on $\fO_Y \oplus \fL$ by $\langle a, b \rangle \cdot \langle a', b' \rangle = \langle aa' + \varphi(b \otimes b'), ab' + a'b \rangle$, where $\varphi$ is an isomorphism of $\fL \otimes \fL \to \fO_Y$.
  Then take $X = \rSpec(\fO_Y \oplus \fL)$.
  Show that $X$ is an étale cover of $Y$.
  \item[(c)] Show that these two processes are inverse to each other.
\end{itemize}

\begin{proof} $ $ \vspace{0pt}
\begin{itemize} [leftmargin=0cm]
  \item[(b)] For each point $y \in Y$, we have $X_y \cong \spec(\fO_y^{\oplus 2})$ with $\fO_{y}$-algebra structure given by
  \begin{equation*}
    \langle a, b \rangle \cdot \langle a', b' \rangle = \langle aa' + bb', ab' + a'b \rangle.
  \end{equation*}
  In particular, the natural morphism $f : X \to Y$ is flat because $f_*\fO_X$ is locally free of rank 2, and the structure morphism $\fO_{y} \to \fO_{y}^{\oplus 2}$ is given by $1 \mapsto \langle 1, 0 \rangle$.

  \item[(c)]
\end{itemize} 
\end{proof}


\end{enumerate}

\end{document}